
\begin{proposition}[A locally unibranch tangent-root component gives positive Euler number]\label{cdp:st:C.2.RAT.5}
Retain the torus-fibre, rational-normalization and specified
ruled-boundary morphism of \ref{cdp:st:C.2.RAT.1}. Suppose $C$ is not
contracted and maps with degree $d=2$ or $3$ onto a smooth
rational curve $H$ with generic ordinary-triple transversal.
Let $q_3\in H^0(H,\operatorname{Sym}^3N^*\otimes M)$ be
its cubic normal tensor, where $N=N_{H/Y}$ and
$M=\OO_Y(S)|_H$. Suppose $q_3$ is nonzero as a tensor at
every point of $H$. Its tangent-root divisor
$R_3\subset\PP(N)$ is then finite of degree three over $H$.
Let $R_d$ be the component corresponding to $C$ over the
generic point. If $R_d$ is locally unibranch at every point,
then
\[
                           e(S)\ge d-1>0.
\tag{C.RD.1}\label{cdp:eq:C.RD.1}
\]
In particular this applies when $d=2$ and the only singularities
of $R_2$ are an $A_4$ cusp or two $A_2$ cusps.
\end{proposition}
\begin{proof}
The nonzero cubic has three distinct roots over the generic
point and is never identically zero on a projective fibre.
It therefore defines a reduced finite divisor $R_3$.
The function-field algebra of its roots agrees with that of
the normalized conductor curves dominating $H$: a branch
of an ordinary triple is distinguished by its tangent line.
Taking the integral closure of $\OO_H$ in that algebra gives
the exact finite-cover identification
\[
                            C\simeq R_d^\nu
                            \quad\text{over }H.
\tag{C.RD.2}\label{cdp:eq:C.RD.2}
\]
Thus two distinct points of $C$ over one point $p\in H$
give distinct root points of $R_d$, because normalization
of a locally unibranch reduced curve is bijective on its
underlying points.

We check what a simple root implies for the actual
hypersurface, including its conductor. In coordinates
$H=(x,y)=0$, with parameter $z$ on $H$, the local equation
$f_0$ of $S$ belongs to $(x,y)^3$. The lower normal jets
vanish on the reduced smooth curve, and the graded pieces
of its regular-sequence ideal are locally free, which
justifies this membership. Choose the normal coordinates
so a specified simple tangent root is finite in $x/y$.
Then
\[
 f_0(y\lambda,y,z)/y^3
       =q_3(\lambda,1,z)+y\,A(\lambda,y,z).
\tag{C.RD.3}\label{cdp:eq:C.RD.3}
\]
The simple root gives a holomorphic solution
$\lambda=\lambda(y,z)$ by the implicit-function theorem.
Division by the monic linear factor yields an actual
factorization
\[
 f_0=(x-y\lambda(y,z))\,G(x,y,z),\qquad
 G|_{x=y\lambda(y,z)}=y^2 u(y,z),\quad u(0)\ne0.
\tag{C.RD.4}\label{cdp:eq:C.RD.4}
\]
The unit is the derivative of the cubic with respect to
its root coordinate, at the chosen simple root. Hence
this normalization branch is smooth and its conductor
in $S$ is exactly twice the curve lying over $H$.
There is no vertical conductor component there, and
the minimal resolution $\pi$ is an isomorphism nearby.

Two distinct $C$-points over $p$ cannot both receive a
vertical boundary attachment. Their root points are
distinct by \eqref{cdp:eq:C.RD.2}. If both could receive attachments,
\eqref{cdp:eq:C.RD.4} would make both roots multiple, requiring total
cubic degree at least four. Thus each point of $H$
receives at most one attachment from $C$.

There are at most $3-d$ generic normalization branches
left to be supplied by vertical strict conductor
components. In the notation of \ref{cdp:st:C.2.RAT.3} this gives
$r_{\rm dom}\le3-d$. At least
$K\ge r-N_0-r_{\rm dom}$ original positive vertical
boundaries therefore survive with an image different
from $H$. Their $K$ points of $C$ give $K$ distinct
normalization-incidence edges at $H$, by the preceding
paragraph. Equation (C.RA.24) gives $n_G\ge K$, and
the exact Euler decomposition gives
\[
 e(S)\ge2-r+N_0+K
       \ge2-r_{\rm dom}\ge d-1.
\]
This proves the assertion. The final examples are locally
unibranch by their displayed cusp equations.
\end{proof}


